\newcommand{\ModuleID}{EL-II.13}
\newcommand{\ModuleTitle}{Place, Time, Means, Manner, and Cause}
\newcommand{\ModuleStage}{Stage II — Core Grammar}
\newcommand{\ModuleUnit}{II.13}
\newcommand{\ModuleOutcome}{Rank morphological and semantic evidence to distinguish spatial, temporal, instrumental, agentive, associative, manner, source, and causal case uses}
\newcommand{\ModulePrerequisites}{EL-II.12 mastery: distinguish verbal nouns/adjectives and obligation structures}
\newcommand{\ModuleExtensionPractice}{Worksheets II.50 and II.51 remain optional remediation or delayed-recall variants}
\newcommand{\ModuleNext}{EL-II.14, Word Order, Ellipsis, Attraction, and Liturgical Compression}
\newcommand{\ModuleMastery}{Label and translate at least 18 of 20 mixed adverbial phrases, produce all six taught exceptional place forms, and compose the ten contrasts without confusing agent, means, location, or direction.}
\newcommand{\ModuleDelayNote}{}
\newcommand{\ModuleLessonSource}{\CourseRoot/shared/content/volume2/all}
\newcommand{\ModuleExerciseSource}{\CourseRoot/shared/exercises/volume2/all}
\newcommand{\ModuleWorkedBank}{II}
\newcommand{\ModuleExerciseBank}{II}
\newcommand{\ModuleWorksheetVariants}{A,D}
\newcommand{\ModuleScopeNote}{This packet gives a ranked operational protocol, not permission to assign a traditional ablative label without syntactic and semantic evidence.}
\newcommand{\ModulePractice}{%
  \SubmoduleExtension
    {The adverbial cases}
    {A ranked protocol prevents the broad label \English{ablative} from standing in for a complete analysis.}
    {For each lesson category, answer in order: exact form and case, governing preposition or verb, animate/personal versus thing distinction, spatial or temporal relation, and narrowest semantic label. Turn the order into a one-page decision card.}
    {Recite the protocol and identify the step that separates personal agent from means.}
    {The protocol begins with morphology, then governor/preposition, lexical type, and contextual relation before assigning the traditional label. Personal agent ordinarily involves a person with \Latin{a/ab} after a passive; means ordinarily uses an unprepositioned thing in the ablative.}
  \SubmoduleExtension
    {Reading the forms in context}
    {Minimal spatial and temporal pairs show how case and preposition alter relation without changing the noun's lexical meaning.}
    {Pair the lesson's location/direction, time-when/duration, agent/means, accompaniment/manner, and source/cause examples. Underline the formal difference and state the semantic difference separately.}
    {Explain why English \English{with} does not prove accompaniment or manner.}
    {Each pair must cite the Latin preposition or bare case and the governing context. English \English{with} can render instrument, accompaniment, or manner; the Latin construction and noun relation decide among them.}
  \SubmoduleExtension
    {Writing laboratory}
    {Composing contrasts tests whether the learner can select a construction before choosing an English preposition.}
    {Complete all ten lesson contrasts. Above each phrase write its case, governor/preposition, and narrow label; then check the six exceptional place forms from memory.}
    {State one contrast where the same case serves two different semantic labels and name the deciding evidence.}
    {All contrasts must preserve location versus direction, time point versus duration, agent versus means, accompaniment versus manner, and source versus cause. The chosen example must identify its different governor, preposition, lexical type, or context rather than claiming the case ending alone decides.}
}
